\documentclass[12pt]{article}

% The preamble file contains all the macros and page setup. I recommend making a folder where you compile your exercises and always call on the preamble.
%%%%%%%%%%%%%%%%%%%%%%%%%%%%%%%%%%%%%%%%%%%%%%%%%%%%%%%%%%%%%%%%%%%%%%
%% all the formatting stuff and packages



\usepackage{geometry}
\geometry{
  top=4cm,
  headheight=70pt,
  headsep=25pt,
  inner=1.25in,
  outer=1.25in
}

\usepackage[utf8]{inputenc}

\usepackage{xcolor,graphicx,hyperref,amsthm,amsmath,amssymb}
\usepackage{algpseudocode}


\usepackage{fancyhdr}
\pagestyle{fancy}

%%% Begin Header Definitions %%%
\fancypagestyle{firststyle}
{
   	\lhead{}
	\rhead{\raisebox{0.2cm}{\includegraphics[width=3.5cm]{resource/hpi_logo.pdf}}}
	\renewcommand{\headrulewidth}{0pt}
}


\newcommand{\makeheader}[1]{
\noindent
\textbf{Algorithmix \hfill Winter term 2019/20 \\ Andreas G\"obel}

\begin{center}
\large{\textbf{\makeHeadline}}\\
\end{center}

\begin{center}
\textit{#1} 
\end{center}


}

\newcommand{\lecturer}[1]{\newcommand{\makeLecturer}{#1}}
\newcommand{\lectureName}[1]{\newcommand{\makeLectureName}{#1}}
\newcommand{\semester}[1]{\newcommand{\makeSemester}{#1}}
\newcommand{\headline}[1]{\newcommand{\makeHeadline}{#1}}
\newcommand{\urlLine}[1]{\newcommand{\makeUrlLine}{#1}}
%%% End Header Definitions %%%

\newcommand{\pr}{\ensuremath{\mathbb{P}}}
\newcommand{\E}{\ensuremath{\mathbb{E}}}


\newcommand{\makeCompact}{%
\setlength{\itemsep}{1pt}%
\setlength{\parsep}{0pt}%
\setlength{\parskip}{0pt}}


\newcommand{\Var}{\mathrm{Var}}


\newcommand{\student}[2]{
\thispagestyle{firststyle}

%%% Begin Header Specification %%%

\headline{Solutions to homework #1}

\makeheader{#2}

\medskip

}

\newcommand{\sol}[1]{ \hspace{.5cm}\\ \noindent \textbf{Exercise #1.}}

% Here you can define your own macros
\newcommand{\calF}{\ensuremath{\mathcal{F}}}

\begin{document}
 
 %To create a header for the document use the following macro. \student{Homework number}{Name}

 \student{1}{Stefan Reschke}
 

 \sol{1} In the lecture we learned about euclidean algorithm to compute the greatest common divisor (\texttt{gcd}) of two numbers $x,y\in \mathbb{N}$. The least common multiple (\texttt{lcm}) of $x,y$ and the gcd can be put together using the following equation \begin{equation}
 	\texttt{lcm}(x,y)=\frac{x\cdot y}{\texttt{gcd}(x,y)}.
 \end{equation}
 
 \textit{Proof.} First we define $d$ as the \texttt{gcd} of $x,y$ \begin{equation}
 	d:=\texttt{gcd}(x,y)
 \end{equation} and $l$ as the \texttt{lcm} of $x$ and $y$ \begin{equation}
 	l:=\texttt{lcm}(x,y).
 \end{equation}
 The equation states that \begin{equation}
 	l=\frac{xy}{d}.
 \end{equation}
 $\frac{xy}{d}$ obviously is a common multiple of $x$ and $y$ because \begin{equation}
 	\frac{xy}{d}\mid x\quad\textrm{and}\quad\frac{xy}{d}\mid x.
 \end{equation}
 For $\frac{xy}{d}$ to be $l$, the following expression needs to hold: \begin{equation}
 	\lnot\bigg(\exists g<\frac{xy}{d}:\big(x|g\land y|g\big)\bigg)
 \end{equation} This can be simplified as follows: \begin{align}
 	\lnot\bigg(\exists g<\frac{xy}{d}&:\big(x\mid g\land y\mid g\big)\bigg) \\
 	\forall g<\frac{xy}{d}&:\lnot\big(x\mid g\land y\mid g\big) \\
 	\forall g<\frac{xy}{d}&:\big(x\nmid g\lor y\nmid g\big)
 \end{align}
 Let $g$ be a number less than $\frac{xy}{d}$. We can write \begin{align}
 	g&<\frac{xy}{d} \\
 	g\cdot d&<x\cdot y.
 \end{align}
 For the first case, we assume that $x\nmid g$. We know that \begin{equation}
 	\exists k:\ g=k \cdot x
 \end{equation}
 
% In the lecture we discussed that the \texttt{gcd} of $x$ and $y$ is always a linear combination of $x, y$. \begin{equation}
% 	\exists a,b:\ d=ax+by
% \end{equation}
 
 
 \sol{2} (a) Let $N$ be an $n$-Bit long number. In order to see how many Bit we need to store $N!$, we have to calculate the binary logarithm of $N!$. \begin{equation}
 	\log_2(N!) = \log_2\bigg(\prod_{i=1}^{N}i\bigg) = \sum_{i=1}^N \log_2(i)
 \end{equation}
 
 \sol{3}
 \noindent (a) To take a look into this lemma, we test for $p=5$. \begin{equation}
 	1^2\equiv 1,\ 2^2\equiv 4,\ 3^2\equiv 4,\ 4^2\equiv 1 \pmod{5}
 \end{equation}
 Apparently the values at the edge of the given set $\{1, \ldots, p-1\}$ have themself as their modular inverse. \begin{enumerate}
 	\item It is obvious that \begin{equation}
 		1^2 \equiv 1 \pmod{p}
 	\end{equation} holds for every $p$, not only primes.
 	\item We take a look at $p-1$. We see that \begin{equation}
 		(p-1)^2=p^2-2p+1
 	\end{equation} and realize that \begin{equation}
 		p^2-2p \equiv 0 \pmod{p}.
 	\end{equation} Hence, we know that \begin{equation}
 		(p-1)^2 \equiv 1 \pmod{p}.
 	\end{equation}
 \end{enumerate}

 \noindent (b) We can split up the factorial of $p-1$ to \begin{gather}
 	(p-1)!=1\cdot 2\cdot\ldots\cdot (p-2)\cdot (p-1).
 \end{gather} Since every prime $p$ that fulfils $p>2$ is a odd number, we can build pairs from elements of the set \begin{equation}
 	P := \{x\in\mathbb{N}| x<p\}
 \end{equation} because we know that \begin{equation}
 	\exists k\in\mathbb{N}: \# P=2k.
 \end{equation} Since we know that \begin{equation}
 	\forall k_1\in P:\ \exists! k_2\in P:\ k_1\cdot k_2 \equiv 1 \pmod{p}
 \end{equation} $(p-1)!$ shortens itself down very easily. Only $p-1$ (which is its own modular inverse) stays. We see that \begin{equation}
 	(p-1)!\equiv p-1 \equiv -1 \pmod{p},
 \end{equation} which was assumed.
 
 \noindent (c) If $N$ is not prime, we know that \begin{equation}
 	N = \prod_{i\in P} i
 \end{equation} where $P$ is the list of prime factors of $N$. We also know that \begin{equation}
 	\forall p\in P:\ N \equiv 0 \pmod{p}
 \end{equation} holds, since $p \div 
 
 \sol{4}
 
 \sol{5}
 
\end{document}
